\section{平均値差の検定}
\label{D1_lesson19}

\subsection{授業内容}
\subsubsection{科目の中でのこのコマの位置づけ}
\begin{tabular}[t]{cccccc}
    記述統計[4] & 確率[6] & モデル[10] & 推測・推定[2] & 意思決定[3] & 実践[2]
\end{tabular}

\subsubsection{概要}
心理学の実践で最もよく用いられるのが，平均値差の検定である。実験計画のところで学んだように，RCTによって結果としてみられる群平均の差が効果の大きさと考えられることが，この手法の心理学的重要性を高めている。まずは二群の平均値差の検定を考える。対応がある二群の場合は実はひとつの「差の分布」についての判断になるため，これまでと同様の議論ができる。対応がない二群の場合は，線形モデルの最も単純な場合であり，前期で学んだ様々な要因計画は，分散分析という別名で広く知られている。このように，線形モデルに確率的判断の手続きを加えたものが，平均値差の検定であると言える。実践的側面で注意するべき点は，対抗するモデルである帰無仮説がどのようにおかれているかと，モデル比較の結果が出た後それをどのように解釈するか，である。
\subsubsection{コマ主題細目(箇条書き)}
\begin{description}



\item[結果の報告]相関係数の有意性検定について解説する。別名「無相関検定」ともよばれているように，帰無仮説が想定する世界は母相関が0の状態である。ここから得られた標本が作る標本相関の分布は$t$分布に従うことがわかっており，これを使って検定を行うことができる。計算のプロセスを簡単にフォローするとともに，結果の報告における表記法について理解する。
\begin{flushright}
    $\to$\citet{Yamada2004}のP.132--133,\citet{山内201001}のP.216--220.
\end{flushright}

\item[統計的に有意，とは]相関係数の帰無仮説検定においては，顕著な違いが示されること＝相関係数が大きいこと，を意味するわけではない。例えばサンプルサイズが大きくなれば，非常に小さな相関係数であっても統計的に有意であると判断される。有意かどうかにこだわるのではなく，その大きさに基づいて判断することの重要性を学び，標準化された効果の大きさを表現する効果量について理解する。
\begin{flushright}
    $\to$\citet{Yamada2004}のP.142--143.効果量については\citet{大久保201201}のP.43--46.
\end{flushright}


\item[検定における2種類の間違い]ここで検定における2種類の間違い方，すなわちType I ErrorとType II Errorについて解説する。モデルの定量的な比較でなく，仮定に基づく質的な確率的な判断であり，その組み合わせから4つの結果が起こり得る。ここでNHSTでは何をコントロールし，何を目的としているかを改めて確認すること，また検定を繰り返すことでType I Errorがインフレーションを起こしうることを理解することが必要である。
\begin{flushright}
    $\to$ \citet{Yamada2004}のP.120--121，検定の多重性については同書P.158--161.\citet{川端201408}の P.90--93.も
\end{flushright}


\end{description}
\subsubsection{キーワード}
\begin{itemize}
\item 信頼区間
\item Type I/II Error
\item 効果量
\end{itemize}
\subsection{授業情報}
\paragraph{コマの展開方法}
講義
\paragraph{標準シラバスにおける位置づけ}
科目番号5；心理学統計法；1.心理学で用いられる統計手法；H 推測統計:代表値と散布度をめぐって(2)
\subsubsection{予習・復習課題}
\paragraph{予習}
% 前期に学んだ要因計画法は，$F$分布と合わせて分散分析としてよく知られているが，2水準の場合は特に$t$分布を使うため$t$検定とも呼ばれている。これらをキーワードに，事前に類書を参照しておくと良い。
統計的帰無仮説検定の手続きを一通り復習し，特に帰無仮説と対立仮説がどのように用意されているか，意思決定の判断基準はどういう意味であったかをしっかりと復習して挑む。
\paragraph{復習}
検定の話に入ると，要因計画や線形モデルとどのようにこれらが接合するのかが分かりにくい。回帰モデルの説明変数を離散化したのが要因計画$\to$要因計画の効果の大きさ(傾きの大きさ)を質的に判断する手続きが，平均値差の検定である，という関係であることを整理しておこう。推測統計学は，基本的に母数の推定(回帰係数のパラメータも母数である)をしているのであり，推定量を何らかの基準で判断するのが検定，という繋がりである。もちろん検定を行わずに，大きさをそのまま報告しても良い。むしろ昨今ではそちらの方が推奨されており，些細な違いや「差がないこと」を「差があること」の根拠とすることの問題点が数多く指摘されている。帰無仮説検定の手続きは，過去の論文に示された結果を読むための知識として重要であり，本質である効果の大きさを推定していることを忘れないように。